% !TEX root = ../main.tex
\chapter{Trabajo Futuro}
\label{cap:futuro}

Los resultados obtenidos en este trabajo abren múltiples líneas de investigación y desarrollo que no han podido abordarse en el marco temporal y de recursos del TFM. Este capítulo describe de forma detallada las limitaciones identificadas y las extensiones más prometedoras.

\section{Ejecución en hardware cuántico real}

La limitación más significativa de este trabajo es que todos los experimentos se han realizado en simulación clásica mediante el simulador de vectores de estado de Qiskit. Si bien la simulación exacta garantiza resultados sin ruido y es reproducible, no refleja las condiciones reales del hardware cuántico actual. Las siguientes extensiones son prioritarias:

\begin{itemize}
    \item \textbf{Ejecución en IBM Quantum (procesadores Eagle/Osprey/Heron):} los notebooks están diseñados para ser ejecutados directamente en hardware real mediante el cambio del backend de simulación a un procesador IBM Quantum. El principal desafío será el ruido de decoherencia, que degradará la calidad del kernel cuántico. Se necesitará implementar técnicas de mitigación de errores (\textit{error mitigation}) como ZNE (\textit{Zero-Noise Extrapolation}) o técnicas de regularización del kernel \citep{haug2023}.
    \item \textbf{Análisis de la degradación del rendimiento con ruido:} comparar sistemáticamente los resultados obtenidos con el simulador ideal frente al hardware ruidoso para cuantificar la reducción de ventaja cuántica debida al ruido.
    \item \textbf{Acceso a hardware cuántico de SKG:} si la empresa obtiene acceso a servicios de computación cuántica en la nube (IBM Quantum, Azure Quantum, Amazon Braket), sería posible ejecutar los experimentos de optimización QAOA sobre los datos reales de movilidad directamente en hardware.
\end{itemize}

\section{Escalabilidad del kernel cuántico}

El coste $O(n^2)$ de la evaluación del kernel cuántico es el principal obstáculo para su aplicación a datasets grandes. Las siguientes líneas permitirían abordar este problema:

\begin{itemize}
    \item \textbf{Kernels cuánticos aleatorizados (RFF-QK):} aplicar la técnica de características aleatorias de Fourier (\textit{Random Fourier Features}) adaptada al contexto cuántico \citep{haug2023}, que permite aproximar el kernel cuántico en tiempo $O(n)$.
    \item \textbf{Particionado y agregación (\textit{ensemble} cuántico):} entrenar múltiples QSVM sobre subconjuntos del dataset y combinar sus predicciones, reduciendo el tamaño efectivo del problema de kernel.
    \item \textbf{Caché y precómputo de kernels:} para experimentos repetitivos, el almacenamiento de las matrices de kernel evaluadas en disco permitiría reutilizar los cálculos sin necesidad de re-ejecutar las simulaciones cuánticas.
\end{itemize}

\section{Ampliación de casos de uso en SKG Tecnología}

Los datasets de SKG Tecnología utilizados en este trabajo son una fracción de la información disponible en el sistema de gestión de movilidad urbana. Las siguientes extensiones son directamente aplicables:

\begin{itemize}
    \item \textbf{Clasificación multiclase de gravedad de incidentes:} ampliar el problema binario actual a una clasificación de tres clases (simple, con lesionado, fatal), lo que requeriría estrategias de kernel multiclase (\textit{one-vs-rest} o \textit{one-vs-one}).
    \item \textbf{Optimización de rutas de respuesta a emergencias:} formular el problema de asignación de vehículos de emergencia como un TSP con restricciones de tiempo (VRP-TW) y resolverlo con SA-QUBO, usando los datos geográficos reales de Bogotá (\texttt{bctraveltimes\_geom.csv}).
    \item \textbf{Predicción de nivel de servicio (clasificación temporal):} usar los datos de tiempos de viaje para predecir el nivel de servicio futuro (verde/amarillo/rojo) como un problema de clasificación con features temporales, evaluando si el kernel cuántico captura patrones de congestión recurrente.
\end{itemize}

\section{Escalación del patrón híbrido al problema real de asignación de SKG}

El Caso de uso~7 demostró el patrón QAOA-COBYLA sobre una red social de 6 nodos con formulación de conjunto dominante mínimo. El paso natural siguiente es escalar este mismo patrón al \textbf{problema real de asignación de agentes a focos de seguridad} de SKG Tecnología, que ya está parcialmente caracterizado en el Caso~3 y cuyos datos están disponibles en \texttt{agents.csv}.

La instancia real presenta una estructura más compleja que el MDS: los agentes tienen capacidades heterogéneas, los focos tienen niveles de criticidad variables y las distancias son métricas geoespaciales (Haversine). Sin embargo, el mismo pipeline híbrido es aplicable con las siguientes adaptaciones:

\begin{itemize}
    \item \textbf{Formulación QUBO de asignación ponderada:} la función de coste incorpora la distancia Haversine entre agente y foco, ponderada por criticidad del incidente. La formulación QUBO resultante tiene estructura similar al \textit{Weighted Set Cover}, que es también NP-difícil \citep{garey1979} y se mapea directamente a un Hamiltoniano de Ising.

    \item \textbf{Validación cruzada sobre \texttt{agents.csv}:} el dataset contiene 10--20 agentes y 5--10 zonas geoespaciales activas, lo que equivale a $n \in [50, 200]$ variables binarias en la codificación QUBO estándar (una por par agente-zona). Esta escala es viable en hardware IBM Eagle (127 qubits) con $p = 1$ y representa la primera instancia real de QAOA sobre datos operativos de SKG.

    \item \textbf{Comparativa con el solver clásico existente:} el módulo \texttt{optimizacion-recursos-campo} ya implementa el solver clásico PuLP/SCIP. La comparativa directa QAOA vs ILP sobre los mismos datos reales completaría la evidencia empírica del TFM con un caso de uso operativamente relevante.

    \item \textbf{Protocolo de activación selectiva:} definir un umbral de activación —por ejemplo, número de agentes disponibles, nivel de congestión del sistema o urgencia del incidente— a partir del cual el módulo QAOA se activa en lugar del greedy, implementando así el criterio de la Tabla~\ref{tab:flujo_hibrido} del Capítulo~\ref{cap:metodologia} sobre datos reales.
\end{itemize}

Esta extensión convertiría el prototipo híbrido descrito en la Sección~\ref{sec:prototipo_hibrido} en una implementación ejecutable sobre infraestructura real, cerrando el ciclo entre la investigación académica y la aplicación operativa en SKG Tecnología.

\section{Mejora de los algoritmos variacionales}

Los experimentos con VQC y QAOA utilizaron configuraciones conservadoras ($p=1$, \texttt{reps=2}) para garantizar la ejecución en simulación. Las siguientes mejoras podrían incrementar la ventaja cuántica:

\begin{itemize}
    \item \textbf{QAOA con mayor profundidad ($p > 1$):} estudios previos \citep{zhou2020} demuestran que la calidad de la solución mejora monotónamente con $p$. Evaluar $p \in \{2, 3, 4\}$ sobre las instancias de optimización de SKG.
    \item \textbf{Optimizadores avanzados:} sustituir COBYLA por optimizadores cuánticos más eficientes como SPSA (\textit{Simultaneous Perturbation Stochastic Approximation}) o métodos de gradiente analítico (\textit{parameter-shift rule}) para acelerar la convergencia.
    \item \textbf{Ansätze específicos del problema:} diseñar circuitos variacionales adaptados a la estructura del problema de optimización (por ejemplo, ansätze con simetría de permutación para problemas de asignación).
\end{itemize}

\section{Integración en pipelines productivos}

El prototipo híbrido cuántico-clásico conceptual diseñado en este trabajo requiere trabajo de ingeniería adicional para su integración real en los sistemas de SKG:

\begin{itemize}
    \item \textbf{API de inferencia cuántica:} desarrollar un microservicio REST que encapsule el QSVM entrenado y exponga un endpoint de predicción, integrable en el pipeline de análisis de incidentes existente.
    \item \textbf{Monitorización de la degradación del modelo:} implementar métricas de \textit{data drift} que detecten cuando los datos de entrada se alejan de la distribución de entrenamiento, desactivando el módulo cuántico en favor del clásico cuando la confianza es baja.
    \item \textbf{A/B testing cuántico-clásico:} diseñar un experimento controlado en producción que evalúe el rendimiento del QSVM frente al Gradient Boosting sobre datos reales en tiempo real.
\end{itemize}

\section{Extensión del análisis teórico}

Finalmente, desde una perspectiva académica, las siguientes extensiones teóricas complementarían los resultados empíricos:

\begin{itemize}
    \item \textbf{Certificación de ventaja cuántica:} aplicar los tests estadísticos propuestos por \cite{huang2021} para certificar formalmente si los datasets reales de SKG presentan la propiedad de separación cuántica-clásica necesaria para garantizar ventaja.
    \item \textbf{Análisis de barren plateaus:} estudiar la presencia de gradientes exponencialmente pequeños en los VQC entrenados, evaluando si las configuraciones utilizadas están afectadas por este fenómeno \citep{cerezo2023}.
\end{itemize}

\section[Extensión del score de cuanticidad]{Extensión del score de cuanticidad $\mathcal{Q}(\mathcal{D})$}

El Capítulo~\ref{cap:investigacion_original} propone el score compuesto $\mathcal{Q}(\mathcal{D}) = w_1 \Delta A + w_2 \hat{d}_F + w_3 \Delta H + w_4 \hat{g}$. El barrido NB08 ($n=9$, $X$ fijo) muestra $\rho_S(\alpha,\Delta A)=0.833$ ($p=0.005$), pero \textbf{no} valida aún a $\mathcal{Q}$ como predictor de ventaja bajo un protocolo structure-aware ($d_F$/$H$ quedaron constantes). La validación sobre tres datasets reales ($n=3$) es orientativa. Extensiones necesarias:

\begin{itemize}
    \item \textbf{Validación sobre benchmarks con ventaja garantizada:} evaluar $\mathcal{Q}(\mathcal{D})$ sobre datasets de Bravyi et al. (2020) donde existe separación cuántico-clásica demostrada formalmente, verificando que $\mathcal{Q} > 0$ en todos ellos.
    \item \textbf{Optimización de los pesos $(w_1, w_2, w_3, w_4)$:} los pesos actuales $(0.4, 0.2, 0.2, 0.2)$ son una convención heurística. Una regresión supervisada sobre un corpus de datasets anotados con su ventaja QSVM real permitiría calibrarlos de forma óptima.
    \item \textbf{Extensión a más feature maps:} el módulo \texttt{feature\_map\_alignment} ya contrastó ZZ vs Pauli-Y en sintéticos. Queda escalar a \texttt{IQP Encoding}, \texttt{TwoLocal} y datasets reales, recalibrando $\mathcal{Q}$ por familia de circuito.
    \item \textbf{Certificación estadística de ventaja:} aplicar los tests estadísticos de \cite{huang2021} para certificar formalmente si los datasets reales de SKG presentan la propiedad de separación cuántica-clásica necesaria para garantizar ventaja, más allá de la estimación empírica del score.
    \item \textbf{Análisis de barren plateaus:} estudiar la presencia de gradientes exponencialmente pequeños en los VQC entrenados, evaluando si las configuraciones utilizadas están afectadas por este fenómeno \citep{cerezo2023}.
\end{itemize}

\paragraph{Perspectivas de publicación.}
La combinación de la propuesta del score $\mathcal{Q}$, su validación empírica (incluyendo el barrido sintético de $n=9$ con $p=0.005$) y el análisis de sensibilidad dimensional constituye el núcleo de un artículo de investigación. El trabajo podría estructurarse como: (i) revisión de las condiciones teóricas de ventaja cuántica en kernels, (ii) propuesta del score como proxy operativo, (iii) validación empírica sistemática sobre corpus anotado, y (iv) primeros pasos hacia la demostración formal. Esta perspectiva fue discutida explícitamente con el director del TFM.

\section{Extensión a aprendizaje cuántico no supervisado}
\label{sec:futuro_no_supervisado}

El Caso de uso~8 (Capítulos~\ref{cap:metodologia} y~\ref{sec:eval_anomaly}) ya incorpora una primera familia \textbf{one-class}: detección de anomalías con OCSVM-ZZ frente a IsolationForest / LOF / OCSVM-RBF, con ventaja cuando la novedad está alineada con el feature map (\texttt{zz\_phase} es la familia que valida ese principio para ZZ). Extender a otros mapas (Pauli, IQP) y a clustering espectral / Q-means / autoencoders temporales queda como trabajo futuro.

\paragraph{Algoritmos candidatos (pendientes).}
\begin{itemize}
  \item \textbf{Q-means \citep{kerenidis2019}:} versión cuántica del algoritmo $k$-means que aprovecha la representación en amplitudes para calcular centroides y distancias en tiempo sublineal. Aplicable al clustering de incidentes de tráfico por zona geoespacial o por patrón horario, un problema que SKG Tecnología aborda actualmente con $k$-means clásico. El módulo QML~v2 \texttt{clustering.py} (Spectral-ZZ) es un primer paso empírico; Q-means con qRAM sigue siendo futuro.

  \item \textbf{Quantum PCA (qPCA) \citep{lloyd2014}:} extracción de componentes principales mediante estimación de fase cuántica aplicada al operador de densidad del dataset. Relevante para la reducción dimensional de los datasets de movilidad antes de aplicar el QSVM, potencialmente superando la PCA clásica en datasets con estructura cuántica implícita.

  \item \textbf{Quantum autoencoders \citep{romero2017}:} circuitos variacionales entrenados para comprimir y reconstruir estados cuánticos. Complementarían el Caso~8 (OCSVM-ZZ) en series temporales: patrones de congestión atípica como estados difícilmente reconstruibles.

  \item \textbf{Clustering cuántico espectral a escala:} el Spectral-ZZ del protocolo QML~v2 ya se midió; falta escalar afinidades y contrastar con Q-means / VQE-$k$-means en $N$ operativos.
\end{itemize}

\paragraph{Casos de uso directos en SKG Tecnología.}
\begin{itemize}
  \item \textbf{Clustering de incidentes:} agrupar los 27.326 incidentes del NUSE Bogotá por patrón de características (hora, zona, tipo) sin etiquetas previas. El Spectral-ZZ QML~v2 mostró ARI débil; Q-means / VQE-$k$-means son la siguiente referencia.
  \item \textbf{Detección de anomalías en velocidades:} extender el Caso~8 (hoy spikes inyectados / novedad alineada al mapa en \texttt{zz\_phase}) a \texttt{bctraveltimes\_history.csv} con autoencoders o scores temporales, sin convertir el problema en fraude etiquetado.
  \item \textbf{Segmentación no supervisada de agentes:} usar clustering cuántico sobre \texttt{agents.csv} para identificar perfiles de agente latentes, mejorando la asignación a zonas de seguridad.
\end{itemize}

\paragraph{Limitaciones actuales a superar.}
Las implementaciones NISQ de algoritmos cuánticos no supervisados son más experimentales que las supervisadas: Q-means requiere acceso eficiente a la memoria RAM cuántica (qRAM), cuya implementación física sigue siendo un problema abierto \citep{biamonte2017}; Quantum PCA basado en estimación de fase requiere circuitos profundos incompatibles con los dispositivos NISQ actuales. Por ello, en el corto plazo la línea más viable es la comparativa de \textbf{clustering cuántico variacional} (VQE aplicado a la minimización de la función de coste de $k$-means) frente al $k$-means clásico, siguiendo el mismo protocolo de tres datasets (cuántico-nativo, SKG real, dataset público de referencia) establecido en este TFM.

\paragraph{Ventaja teórica en complejidad.}
Uno de los argumentos más sólidos a favor del aprendizaje no supervisado cuántico es la reducción de complejidad computacional que ofrecen los algoritmos más maduros. Q-means \citep{kerenidis2019} tiene complejidad $O(\text{poly}(\log N, k, d, 1/\varepsilon))$, exponencialmente mejor que el $O(Nkd)$ del $k$-means clásico en función del número de muestras $N$, número de clusters $k$ y dimensionalidad $d$. Quantum PCA \citep{lloyd2014} calcula los vectores propios dominantes del operador de densidad en tiempo $O(\log N)$, frente al $O(N d^2)$ de la PCA clásica. Esta reducción exponencial de complejidad es particularmente relevante para los datasets de movilidad de SKG Tecnología, donde $N > 80.000$ registros y la PCA clásica ya constituye un cuello de botella en pipelines de tiempo real. Con hardware cuántico tolerante a fallos, estos algoritmos podrían transformar radicalmente la viabilidad del aprendizaje no supervisado a escala operativa.

\section{Extensión del Framework QCR}

El Framework QCR (Quantum-Compatible Representation) se \textbf{plantea formalmente} en el Capítulo~\ref{cap:investigacion_original} (definición, protocolo, Conjetura~\ref{conj:qcr}) con evidencia preliminar: Notebook~09 (re-uploading $reps=3$ elevaba $\mathcal{Q}$ en Telco \emph{vs.\ RBF}) y revisión QML~v2 (Sección~\ref{sec:qml_v2_revision}), donde Telco permanece \texttt{classical} y SKG solo muestra un cambio geométrico \texttt{try-ZZ} por rotación, sin lift de accuracy. SKG/Telco no son cuántico-nativos. Este TFM no pretende cerrar la conjetura. Quedan líneas de extensión claras:

% -------------------------------------------------------------------
\subsection{Demostración de la Conjetura del Espacio de Representación Cuántico-Compatible}
\label{sec:futuro_conjetura_qcr}
% -------------------------------------------------------------------

La contribución teórica central del Capítulo~\ref{cap:investigacion_original} es la Conjetura~\ref{conj:qcr}, que propone la existencia de un ciclo completo de inferencia cuántica

\begin{equation}
  \mathbf{x}
  \;\xrightarrow{\;\psi\;}\;
  \mathbf{z}
  \;\xrightarrow{\;f_Q^*\;}\;
  \hat{y}
  \;\Big[\xrightarrow{\;\psi^{-1}\;}\; \mathbf{x}^* \;\text{(solo en optimización)}\Big]
\end{equation}

donde $\psi: \mathcal{X} \rightarrow \mathcal{Z}$ es una transformación ---no necesariamente polinómica, en analogía con la transformada de Fourier, que nace de un kernel--- que convierte cualquier dataset clásico en cuántico-compatible. Esta conjetura permanece abierta, y su demostración o refutación constituye la línea de investigación post-TFM de mayor impacto potencial. La hoja de ruta propuesta es la siguiente:

\begin{enumerate}
  \item \textbf{Caracterización de la clase de distribuciones $p(\mathbf{x},y)$ para las que existe $\psi$.} \\
  Determinar si la existencia de $\psi$ depende de propiedades universales de la distribución (como suavidad, separabilidad de clases o estructura espectral) o si existen distribuciones para las cuales ninguna $\psi$ produce $\mathcal{Q}(\psi(\mathcal{D})) > 0$. Un contraejemplo de esta segunda hipótesis constituiría una refutación de la conjetura.

  \item \textbf{Demostración de invarianza de error: $\varepsilon(f_Q^* \circ \psi) \leq \varepsilon(f_C^*)$.} \\
  Probar que la ventaja cuántica obtenida en $\mathcal{Z}$ no se destruye al componer con $\psi$. Esto requiere una cota de la distorsión informacional introducida por $\psi$ ---análoga a la preservación de separabilidad en PCA de kernels--- en términos del score $\mathcal{Q}$ y de la cota de Huang et al.\ \citep{huang2021}.

  \item \textbf{Back-mapping $\psi^{-1}$ con garantías de factibilidad en optimización.} \\
  Para problemas combinatorios (asignación de agentes, rutas, portafolio), la solución cuántica óptima $\mathbf{z}^* \in \mathcal{Z}$ debe decodificarse en una solución factible $\psi^{-1}(\mathbf{z}^*) \in \mathcal{X}$. Hay que demostrar que $\psi$ admite un decodificador que preserva la factibilidad de las restricciones QUBO bajo la transformación inversa, o diseñar un decodificador aproximado con cota de penalización controlada.

  \item \textbf{Algoritmo de búsqueda de $\psi$ con garantías de aproximación.} \\
  El re-uploading adaptativo encontrado empíricamente en el Notebook~09 es solo una instancia particular de $\psi$ en un espacio de búsqueda mucho mayor. Desarrollar un algoritmo de búsqueda ---posiblemente variacional, tipo quantum metric learning \citep{lloyd2020}--- que encuentre $\psi$ óptima con garantías de convergencia y complejidad polinómica en el tamaño del dataset.
\end{enumerate}

\paragraph{Impacto potencial.}
Si la Conjetura~\ref{conj:qcr} se demostrara en su forma fuerte ---es decir, si existiera $\psi$ para toda distribución--- implicaría que la computación cuántica puede ofrecer ventaja sobre \emph{cualquier} dataset clásico, dado el preprocesamiento adecuado. Aunque el horizonte temporal de esta demostración es incierto, incluso su versión débil ---existencia de $\psi$ para familias específicas de distribuciones--- tendría un impacto metodológico directo en el diseño de pipelines de aprendizaje automático cuántico en entornos empresariales.

\subsection{Extensiones operativas del framework}

\begin{itemize}
    \item \textbf{UMAP cuántico con métrica real:} la Estrategia~4 (Q-UMAP) se implementó con MDS como fallback por ausencia de \texttt{umap-learn}. La implementación con UMAP y distancia cuántica genuina ($d_Q = 1 - K_Q$) puede cambiar cualitativamente los resultados, dado que MDS preserva distancias globales mientras UMAP preserva estructura local.
    \item \textbf{Escalar a 4--8 qubits:} el framework actual opera con $n_q = 2$, forzando compresión PCA que descarta hasta el 69\% de la varianza. Con 4--8 qubits y \texttt{SelectKBest} o selección supervisada, se captura más estructura original. Sin embargo, el análisis de sensibilidad (Tabla~\ref{tab:sensibilidad_dim}) advierte que features muy discriminativas pueden reducir $\mathcal{Q}$: el diseño de la selección de features deberá equilibrar relevancia supervisada y compatibilidad con el feature map.
    \item \textbf{Validación cruzada estratificada:} los resultados actuales son sobre $n=80$ con una sola partición 70/30. Repetir con $k$-fold estratificado ($k=5$) permitiría medir la varianza de los indicadores y construir intervalos de confianza para $\mathcal{Q}(\mathcal{D})$ y la ventaja QSVM.
    \item \textbf{Extensión a estrategias de pre-mapeo no lineales:} el framework actual solo incluye transformaciones lineales (rotación) y de feature map (re-uploading). Incorporar \textit{quantum metric learning} \citep{lloyd2020} ---que aprende el mapa de codificación directamente de los datos--- como quinta estrategia completaría el espacio de búsqueda de $T^*$.
    \item \textbf{Aprendizaje del mapeo óptimo \textit{end-to-end}:} formular la búsqueda de $T^*$ como un problema de optimización donde la función objetivo maximiza $\mathcal{Q}(T(\mathcal{D}))$ con restricción de preservación de información discriminativa (e.g., KL-divergencia acotada entre distribuciones de clases). Esto conecta con \textit{quantum-aware representation learning} \citep{huang2021}.
\end{itemize}

\paragraph{Perspectivas de publicación.}
La combinación del score $\mathcal{Q}(\mathcal{D})$ y el Framework QCR --- validados sobre datos empresariales reales y con código reproducible en Qiskit 2.x --- constituye un pipeline metodológico completo con el perfil de contribución necesario para una publicación en \textit{npj Quantum Information} o \textit{Quantum Machine Intelligence}. Esta perspectiva fue contemplada en la reunión de seguimiento con el director del TFM.

\begin{figure}[!htbp]
\centering
\begin{tikzpicture}[font=\small, thick,
    item/.style={rectangle, rounded corners=4pt, very thick,
                 text width=3.4cm, align=left, inner sep=6pt,
                 minimum height=0.85cm},
    header/.style={rectangle, rounded corners=5pt, very thick,
                   minimum width=3.8cm, minimum height=0.8cm,
                   align=center, font=\footnotesize\bfseries}]

  % ── Cabeceras de horizonte ──
  \node[header, draw=blue!55!black,   fill=blue!10]   (h1) at (0,    0) {Corto plazo\\1--2 años};
  \node[header, draw=orange!55!black, fill=orange!8]  (h2) at (5.0,  0) {Medio plazo\\2--4 años};
  \node[header, draw=violet!55!black, fill=violet!8]  (h3) at (10.0, 0) {Largo plazo\\4+ años};

  \draw[gray!35, thick] (-1.8, -0.48) -- (11.8, -0.48);

  % ── Columna corto plazo ──
  \node[item, draw=blue!40, fill=blue!5, below=0.35cm of h1] (b1)
    {\textbf{Hardware real IBM}\\\scriptsize Qiskit Runtime + EstimatorV2};
  \node[item, draw=blue!40, fill=blue!5, below=0.25cm of b1] (b2)
    {\textbf{API REST de inferencia}\\\scriptsize microservicio QSVM};
  \node[item, draw=blue!40, fill=blue!5, below=0.25cm of b2] (b3)
    {\textbf{A/B test Q vs Clásico}\\\scriptsize datos SKG en producción};
  \node[item, draw=blue!40, fill=blue!5, below=0.25cm of b3] (b4)
    {\textbf{QAOA $p > 1$}\\\scriptsize SPSA / parameter-shift};
  \node[item, draw=blue!40, fill=blue!5, below=0.25cm of b4] (b5)
    {\textbf{Aprendizaje no supervisado Q}\\\scriptsize Q-means / Quantum PCA};

  % ── Columna medio plazo ──
  \node[item, draw=orange!45, fill=orange!5, below=0.35cm of h2] (m1)
    {\textbf{Escalar kernel a 4--8 q}\\\scriptsize SelectKBest + ZZFeatureMap};
  \node[item, draw=orange!45, fill=orange!5, below=0.25cm of m1] (m2)
    {\textbf{QAOA en asignación SKG}\\\scriptsize recursos campo + QUBO real};
  \node[item, draw=orange!45, fill=orange!5, below=0.25cm of m2] (m3)
    {\textbf{Validación score $\mathcal{Q}$}\\\scriptsize corpus $>$20 datasets};
  \node[item, draw=orange!45, fill=orange!5, below=0.25cm of m3] (m4)
    {\textbf{Q-UMAP métrica real}\\\scriptsize distancia cuántica $d_Q$};

  % ── Columna largo plazo ──
  \node[item, draw=violet!45, fill=violet!5, below=0.35cm of h3] (l1)
    {\textbf{Corrección de errores}\\\scriptsize hardware tolerante a fallos};
  \node[item, draw=violet!45, fill=violet!5, below=0.25cm of l1] (l2)
    {\textbf{Demostración Conjetura QCR}\\\scriptsize $\exists\,\psi$; back-mapping $\psi^{-1}$};
  \node[item, draw=violet!45, fill=violet!5, below=0.25cm of l2] (l3)
    {\textbf{QCR generalizado}\\\scriptsize PauliMap + IQP + TwoLocal};
  \node[item, draw=violet!45, fill=violet!5, below=0.25cm of l3] (l4)
    {\textbf{Publicación científica}\\\scriptsize npj QI / Quantum ML};
  \node[item, draw=violet!45, fill=violet!5, below=0.25cm of l4] (l5)
    {\textbf{QML end-to-end}\\\scriptsize quantum metric learning};

  % ── Flechas de progresión entre columnas ──
  \draw[->, very thick, gray!50] (2.05, -3.5) -- (2.75, -3.5);
  \draw[->, very thick, gray!50] (7.05, -3.5) -- (7.75, -3.5);

\end{tikzpicture}
\caption{Roadmap de trabajo futuro organizado en tres horizontes temporales. Las líneas de investigación de corto plazo (azul) son directamente ejecutables con el hardware actual; las de medio plazo (naranja) requieren hardware mejorado o mayor volumen de datos; las de largo plazo (violeta) dependen de avances en hardware tolerante a fallos y consolidación teórica, incluyendo la demostración de la Conjetura~\ref{conj:qcr}.}
\label{fig:roadmap_futuro}
\end{figure}

